% Generated draft: requires mathematical and source review.
\documentclass[11pt]{article}
\usepackage[margin=25mm]{geometry}
\usepackage{fontspec}
\setmainfont{Latin Modern Roman}
\usepackage{amsmath,amssymb,mathtools,unicode-math}
\setmathfont{Latin Modern Math}
\usepackage{enumitem,needspace,xurl,hyperref}
\setlength{\parindent}{0pt}
\setlength{\parskip}{5pt}
\setlength{\emergencystretch}{3em}
\newcommand{\N}{\mathbb{N}}
\newcommand{\Z}{\mathbb{Z}}
\newcommand{\Q}{\mathbb{Q}}
\newcommand{\R}{\mathbb{R}}
\newcommand{\C}{\mathbb{C}}
\newcommand{\F}{\mathbb{F}}
\newcommand{\A}{\mathbb{A}}
\newcommand{\PP}{\mathbb{P}}
\newcommand{\E}{\mathbb{E}}
\newcommand{\eps}{\varepsilon}
\newcommand{\dd}{\,\mathrm{d}}
\begin{document}
\section*{UnsolvedMath ID: 700508}
\textbf{Source problem code:} catalogue entry
\subsection{Definitions and mathematical statement}
Let \(G\) be a locally compact topological group. For this draft, local compactness is understood with the usual Hausdorff convention. Interpret \emph{locally nilpotent} algebraically: every subgroup generated by finitely many elements of \(G\) is nilpotent. Here a group \(H\) is nilpotent if its lower central series, defined by \(\gamma_1(H)=H\) and \(\gamma_{i+1}(H)=[\gamma_i(H),H]\), reaches the trivial subgroup after finitely many steps; \([A,B]\) denotes the subgroup generated by commutators \(a^{-1}b^{-1}ab\), with \(a\in A\) and \(b\in B\).

Assume that every proper closed normal subgroup of \(G\) is compact; explicitly,
\[
\forall N\leq G,\qquad
\bigl(N\trianglelefteq G\ \land\ N\text{ is closed}\ \land\ N\ne G\bigr)
\Longrightarrow N\text{ is compact}.
\]
Must there exist a subgroup \(K\leq G\) such that
\[
K\trianglelefteq G,\qquad K\text{ is open in }G,\qquad K\text{ is compact}?
\]
The desired subgroup \(K\) is not required to be proper. All topological properties refer to the given topology on \(G\). The supplied record does not explicitly specify the Hausdorff convention or define local nilpotence; the conventions stated above require confirmation against the source.
\subsection{Short English statement}
Must every locally compact, locally nilpotent group whose proper closed normal subgroups are all compact contain a compact open normal subgroup? The requested subgroup may be the whole group. This draft uses the Hausdorff convention and algebraic local nilpotence, subject to source review.
\subsection{Sources}
\begin{itemize}
\item Local UnsolvedMath registry, permanent ID \(700508\), exact source code \texttt{catalogue entry}, title \emph{Kourovka Notebook Problem 8.59}. The supplied record labels the problem open; this label is not an independently verified mathematical or literature assessment.
\item \emph{The Kourovka Notebook (Unsolved Problems in Group Theory)}, original problem 8.59, 8th issue (1982); supplied edition reference: 21st issue edition (2026), \texttt{arXiv:1401.0300v48}. Attribution: V. M. Poletskikh and I. V. Protasov, proposed in 1982. The supplied collection description names E. I. Khukhro and V. D. Mazurov as editors of Issues 1--21 (1965--2026). These bibliographic details are supplied attribution and have not been independently verified.
\end{itemize}
\end{document}
